\section{Problema}

A utilização do fatiamento de redes permite que os recursos de uma infraestrutura física sejam compartilhados entre diferentes fatias, de acordo com os requisitos de cada serviço ou aplicação. Entretanto, a definição de como esses recursos devem ser distribuídos entre as fatias constitui um desafio, especialmente em ambientes nos quais a demanda varia ao longo do tempo. Nesse cenário, diferentes estratégias podem ser utilizadas para determinar quais recursos serão alocados, quais fatias serão reutilizadas e quando uma nova busca por recursos deve ser realizada.

Diferentes estratégias podem ser adotadas para realizar essa alocação e adaptar as fatias às variações nas necessidades dos usuários ao longo do tempo. \citeonline{gomes2020strategies} comparam estratégias de fatiamento elástico e mostram que a política utilizada afeta métricas como disponibilidade de banda e quantidade de clientes atendidos. Este estudo delimita-se à comparação, sob as mesmas requisições e topologias, do PETIC, de duas buscas aproximadas para seu custo aditivo (ACO e Q-Learning) e do SA, cuja função objetivo inclui equilíbrio de carga.

A questão de pesquisa é: \textbf{como variam a aceitação de demandas, o equilíbrio da utilização dos enlaces, a reutilização de caminhos e o tempo de decisão dessas quatro implementações em diferentes níveis de carga e topologias, quando avaliadas nos mesmos cenários independentes?} São examinadas as seguintes hipóteses: (H1) as diferenças de aceitação entre PETIC, ACO e Q-Learning podem variar por cenário e não se presume uma ordenação entre eles; (H2) a função multiobjetivo do SA favorece maior equilíbrio dos enlaces, potencialmente acompanhado de caminhos mais longos e menor aceitação sob carga alta; (H3) a busca exata por Dijkstra no PETIC requer menos tempo de CPU que as buscas estocásticas avaliadas. As hipóteses são avaliadas por métricas comuns e, para aceitação, por diferenças pareadas por cenário. Como ACO e Q-Learning aproximam o mesmo custo aditivo já minimizado pelo PETIC, este trabalho não testa se superam o ótimo desse custo; investiga como os mecanismos de busca se refletem nos resultados da alocação sequencial.
